\clearpage
\SourceLeaf{154}\pdfbookmark[1]{Sectio quarta.}{sectio-quarta}
\begin{center}
{\large SECTIO QUARTA.}\par\medskip
\emph{Relationes inter locos plures in spatio.}\par\bigskip
\hypertarget{native-154-article110}{}\pdfbookmark[2]{110.}{article-110}110.
\end{center}
Relationes in hac Sectione considerandae ab orbitae indole independentes solique suppositioni innixae erunt, omnia orbitae puncta in eodem plano cum Sole iacere. Placuit autem, hic quasdam simplicissimas tantum attingere, aliasque magis complicatas et speciales ad Librum alterum nobis reservare.

Situs plani orbitae per duos locos corporis coelestis in spatio plene determinatus est, siquidem hi loci non iacent in eadem recta cum Sole. Quare quum duobus potissimum modis locus puncti in spatio assignari possit, duo hinc problemata solvenda se offerunt.

Supponemus primo, duos locos dari per longitudines et latitudines heliocentricas resp. per $\lambda$, $\lambda'$, $\beta$, $\beta'$ designandas: distantiae a Sole in calculum non ingredientur. Tunc si longitudo nodi ascendentis per $\AscendingNode$, inclinatio orbitae ad eclipticam per $i$ denotatur, erit
\[
\begin{aligned}
\operatorname{tang}\beta&=\operatorname{tang}i\sin(\lambda-\AscendingNode)\\
\operatorname{tang}\beta'&=\operatorname{tang}i\sin(\lambda'-\AscendingNode).
\end{aligned}
\]
Determinatio incognitarum $\AscendingNode$, $\operatorname{tang}i$ hic ad problema in art. 78, II consideratum refertur; habemus itaque, ad normam solutionis primae
\[
\begin{aligned}
\operatorname{tang}i\sin(\lambda-\AscendingNode)&=\operatorname{tang}\beta\\
\operatorname{tang}i\cos(\lambda-\AscendingNode)&=\frac{\operatorname{tang}\beta'-\operatorname{tang}\beta\cos(\lambda'-\lambda)}{\sin(\lambda'-\lambda)};
\end{aligned}
\]
\SourceLeaf{155}ad normam solutionis tertiae autem invenimus $\AscendingNode$ per aequationem
\[
\operatorname{tang}(\tfrac12\lambda+\tfrac12\lambda'-\underset{\rule{.45em}{.35pt}}{\AscendingNode})=\frac{\sin(\beta'+\beta)\operatorname{tang}\tfrac12(\lambda'-\lambda)}{\sin(\beta'-\beta)}
\]
utique aliquanto commodius, si anguli $\beta$, $\beta'$ immediate dantur, neque vero per logarithmos tangentium: sed ad determinandum $i$, recurrendum erit ad aliquam formularum
\[
\operatorname{tang}i=\frac{\operatorname{tang}\beta}{\sin(\lambda-\AscendingNode)}=\frac{\operatorname{tang}\beta'}{\sin(\lambda'-\AscendingNode)}\mathpunct{\cdot}
\]
Ceterum ambiguitas in determinatione anguli $\lambda-\AscendingNode$, vel $\tfrac12\lambda+\tfrac12\lambda'-\AscendingNode$ per tangentem suam ita erit decidenda, ut $\operatorname{tang}i$ positiva evadat vel negativa, prout motus ad eclipticam proiectus directus est vel retrogradus: hanc incertitudinem itaque tunc tantum tollere licet, ubi constat, a quanam parte corpus coeleste a loco primo ad secundum pervenerit; quod si ignoraretur, utique impossibile esset, nodum ascendentem a descendente distinguere.

Postquam anguli $\AscendingNode$, $i$ inventi sunt, eruentur argumenta latitudinum $u'$, $u$ per formulas
\[
\operatorname{tang}u=\frac{\operatorname{tang}(\lambda-\AscendingNode)}{\cos i},\qquad
\operatorname{tang}u'=\frac{\operatorname{tang}(\lambda'-\AscendingNode)}{\cos i},
\]
quae in semicirculo primo vel secundo accipienda sunt, prout latitudines respondentes boreales sunt vel australes. His formulis adhuc sequentes adiicimus, e quibus, si placet, una vel altera ad calculum confirmandum in usum vocari poterit:
\[
\begin{aligned}
\cos u&=\cos\beta\cos(\lambda-\AscendingNode),&\quad \cos u'&=\cos\beta'\cos(\lambda'-\AscendingNode),\\[3pt]
\sin u&=\frac{\sin\beta}{\sin i},&\sin u'&=\frac{\sin\beta'}{\sin i},
\end{aligned}
\]
\[
\begin{aligned}
\sin(u'+u)&=\frac{\sin(\lambda+\lambda'-2\AscendingNode)\cos\beta\cos\beta'}{\cos i},\\[3pt]
\sin(u'-u)&=\frac{\sin(\lambda'-\lambda)\cos\beta\cos\beta'}{\cos i}\mathpunct{\cdot}
\end{aligned}
\]
